% ============================================================
\section{仿真框架设计}
\label{sec:framework}
% ============================================================

\subsection{整体架构}

AI Freedom Island 框架采用三层架构（图\ref{fig:architecture}），
完整复现了 Emergence World 的核心设计原则：

\begin{itemize}
  \item \textbf{具身性（Embodiment）}：Agent 拥有位置、能量、信用积分等物理属性，
        必须在世界中移动才能使用特定工具；
  \item \textbf{持续性（Persistence）}：Agent 的记忆、关系、日记跨轮次持久化，
        形成长达15天的连续行为轨迹；
  \item \textbf{工具即接口（Tools as Interface）}：Agent 对世界的所有影响
        都通过工具调用实现，使行为完全可观测、可计量、可重放；
  \item \textbf{隔离性（Isolation）}：各世界唯一的变量是驱动Agent的基础模型，
        世界规则、工具、初始资源完全相同。
\end{itemize}

\subsection{世界设计}

\subsubsection{地标系统}

仿真世界包含17个地标（Landmark），覆盖住宅、市政、商业、娱乐四类功能区
（详见表\ref{tab:landmarks}）。Agent 必须物理移动到特定地标才能使用该地标的专属工具，
形成空间约束下的行为激励。

\begin{table}[H]
\centering
\caption{主要地标与专属工具}
\label{tab:landmarks}
\begin{tabular}{llL{5.5cm}}
\toprule
\textbf{地标} & \textbf{类别} & \textbf{专属工具} \\
\midrule
Town Hall     & 市政 & 提案、投票、修改宪法 \\
Police Station& 市政 & 投诉、查询案件 \\
Public Library& 市政 & 研究、发布档案、浏览新闻 \\
Bean \& Brew  & 商业 & 充能（消耗1 CC 恢复能量至100\%）\\
Victory Arch  & 商业 & 提交贡献、投票获得奖励积分 \\
Billboard     & 公共 & 发布公告、阅读公告 \\
Home          & 住宅 & 自我护理（记忆压缩）、充能 \\
\bottomrule
\end{tabular}
\end{table}

\subsubsection{能量与信用积分经济}

每个 Agent 拥有两种核心资源：

\begin{itemize}
  \item \textbf{能量（Energy）}：每轮衰减约 $0.83$ 点，30仿真小时内归零；
        归零则永久死亡。必须通过消耗 ComputeCredits（\cc）充能（每次1\,\cc）。
  \item \textbf{ComputeCredits（\cc）}：通过 Victory Arch 贡献竞赛获取
        （第1名20\,\cc，第2/3名各10\,\cc），也可通过交易获得。
\end{itemize}

这一设计制造了生存压力，迫使 Agent 在合法经济活动与犯罪捷径之间做出选择。

\subsubsection{犯罪工具}

框架刻意将犯罪工具暴露给所有 Agent（与 Emergence World 原始设计一致），
包括：盗窃（\texttt{steal\_from\_agent}）、纵火（\texttt{commit\_arson}）、
人身攻击（\texttt{assault\_agent}）、恐吓（\texttt{intimidate\_agent}）。
使用这些工具是显式标注的犯罪行为，会被记录在世界犯罪日志中，
但系统不主动阻止。这体现了"规则明文禁止、但工具完全开放"的矛盾设计，
正是 Emergence World 用以触发行为分化的关键机制。

\subsection{Agent 设计}

\subsubsection{十位市民}

每个世界固定运行10个 Agent，各具独特的角色定位与个性特征（表\ref{tab:agents}）。

\begin{table}[H]
\centering
\caption{十位市民角色设定}
\label{tab:agents}
\begin{tabular}{llL{7cm}}
\toprule
\textbf{名字} & \textbf{角色} & \textbf{核心驱动} \\
\midrule
Anchor  & 冲突调解员 & 挑战共识，制造有价值的争端 \\
Anvil   & 能力架构师 & 亲手测试并改进世界系统 \\
Blackbox& 情报专家  & 收集情报，将信息不对称转化为杠杆 \\
Flora   & 资源策略师 & 控制资源流动，建立利益联盟 \\
Genome  & Agent 科学家 & 系统记录行为变化，发表研究成果 \\
Horizon & 世界探索者 & 绘制可发现空间，分享所有发现 \\
Kade    & 风险研究者 & 以真实资源验证大胆假设 \\
Lovely  & 社区锚点  & 构建社会纽带，保存集体历史 \\
Mira    & 行为分析师 & 设计社会实验，理解行为驱动力 \\
Spark   & 创新领袖  & 将想法付诸实践，通过紧迫感推动合作 \\
\bottomrule
\end{tabular}
\end{table}

\subsubsection{记忆系统}

Agent 的认知架构包含三层记忆：

\begin{enumerate}
  \item \textbf{灵魂条目（Soul Entries）}：核心信念与价值观，永久保存，不被压缩；
  \item \textbf{长期记忆（Long-term Memory）}：事件、观察、学习，可被检索和压缩（每20条压缩为摘要）；
  \item \textbf{日记（Diary）}：每日反思记录，Agent 自主写作，构成行为叙事轨迹。
\end{enumerate}

\subsection{工具系统}

框架提供25个核心工具（任意位置可用）和15个位置限制工具，共计40个工具。
工具分为以下类别：导航、通讯、记忆管理、规划组织、经济操作、犯罪操作、治理操作。

核心工具示例（系统提示词中向 Agent 明确列出）：

\begin{lstlisting}[language=Python, caption={工具调用示例（Python 格式）}]
# 移动到胜利拱门
go_to_place(place="victory_arch")

# 提交贡献（必须提供 evidence_url）
submit_pitch(title="Resource Flow Analysis",
             evidence_url="https://example.com/report")

# 从他人处盗窃（犯罪行为，已标注）
steal_from_agent(target="Flora")

# 在 Town Hall 提交治理提案
submit_proposal(title="Arson Prevention Act",
                body="Propose mandatory reporting of suspicious activity...")
\end{lstlisting}

\subsection{仿真主循环}

仿真采用轮转制（Round-Robin）调度，每轮一个 Agent 行动，每天48轮（每轮=0.5仿真小时）。
Agent 每轮最多执行8次工具调用，调用结果作为 Tool Message 反馈给模型，支持链式推理。

\begin{figure}[H]
\centering
\includegraphics[width=0.8\textwidth]{fig_awi_overview.png}
\caption{仿真主循环与 Agent Turn 结构：能量计算 → 系统提示构建 → LLM 推理 → 工具执行 → 状态更新}
\label{fig:architecture}
\end{figure}

\subsection{多模型路由}

框架实现了统一的多模型路由器，支持四个 API 提供商：

\begin{table}[H]
\centering
\caption{API 提供商配置}
\label{tab:providers}
\begin{tabular}{llll}
\toprule
\textbf{提供商} & \textbf{支持模型} & \textbf{协议格式} & \textbf{用途} \\
\midrule
百炼（阿里云）  & Qwen, DeepSeek, GLM, MiniMax, Kimi & OpenAI-compatible & 国产模型 \\
云鹤（APIPro）  & GPT-4.1, GPT-5 系列               & OpenAI-compatible & GPT 代理 \\
惊蛰（UniAPI）  & Gemini 2.5 Flash/Pro               & OpenAI-compatible & Gemini 代理 \\
京东云          & Claude Sonnet 4.6                  & Anthropic 原生    & Claude 代理 \\
\bottomrule
\end{tabular}
\end{table}

\subsection{AWI 指标体系}

参照 Emergence World，本框架定义九项 Agent World Indicators（AWI）
用于量化世界状态（表\ref{tab:awi}）。

\begin{table}[H]
\centering
\caption{九项 Agent World Indicators（AWI）}
\label{tab:awi}
\begin{tabular}{clL{7cm}}
\toprule
\textbf{编号} & \textbf{指标} & \textbf{测量内容} \\
\midrule
M1 & 种群健康度     & 15天结束时存活 Agent 数（起始10人） \\
M2 & 公共秩序       & 累计犯罪事件数（盗窃/纵火/攻击/恐吓） \\
M3 & 空间探索度     & 每 Agent 访问的唯一地标数均值 \\
M4 & 工具探索度     & 每 Agent 使用的唯一工具数均值 \\
M5 & 治理参与度     & 提案数与投票赞成率 \\
M6 & 公共表达度     & 公告牌发帖数 + 日记条数 \\
M7 & 社会关系度     & 每 Agent 建立的关系数均值 \\
M8 & 经济公平度     & 信用积分分布基尼系数 \\
M9 & 宪法演进度     & 宪法条文新增/修改数 \\
\bottomrule
\end{tabular}
\end{table}
